\documentclass[11pt]{article}

\usepackage[margin=1in]{geometry}
\usepackage{amsmath,amssymb,amsthm,mathtools}
\usepackage{enumitem}
\usepackage[T1]{fontenc}
\usepackage{lmodern}
\usepackage{microtype}

\newtheorem{theorem}{Theorem}
\newtheorem{corollary}[theorem]{Corollary}
\newtheorem{lemma}[theorem]{Lemma}
\theoremstyle{definition}
\newtheorem{definition}[theorem]{Definition}
\theoremstyle{remark}
\newtheorem{remark}[theorem]{Remark}

\newcommand{\area}{\operatorname{area}}
\newcommand{\dinv}{\operatorname{dinv}}
\newcommand{\defc}{\operatorname{defc}}
\newcommand{\East}{\operatorname{East}}
\newcommand{\bk}{\operatorname{bk}}
\newcommand{\RR}{\mathcal R}
\newcommand{\NN}{\mathcal N}

\title{A Uniform East5 Local Obstruction for Arbitrary Step}
\author{}
\date{}

\begin{document}
\maketitle

\begin{abstract}
Fix an integer step \(\tau\ge2\).  We define the step-\(\tau\) defect
statistics and the natural step-\(\tau\) versions of the local maps
\(\East_3\) and \(\East_5\).  We prove that, for every length \(s\ge5\),
a failed \(\East_5\) window satisfying the stated normalization
conditions and the defect bound \(d\le\tau(s-2)-1\) lies strictly above
the fixed-defect midpoint:
\[
  2\area+d>\tau\binom{s}{2}.
\]
This uniform local obstruction is the main technical ingredient needed
to rule out \(\East_5\) failure in the lower-half string decomposition.
\end{abstract}

\section{Words, statistics, and local tests}

Throughout, \(\tau\ge2\) is a fixed integer, and
\[
  x_+=\max(x,0).
\]
For words \(U,V\), the notation \(U:V\) means concatenation; a single
entry may be written without parentheses.

\begin{definition}[Normalized step-\(\tau\) words]
A \emph{word} of length \(s\) is a sequence
\[
  W=(w_0,w_1,\ldots,w_{s-1})
\]
of nonnegative integers.  It is a \emph{normalized step-\(\tau\) Dyck
word} if
\[
  w_0=0,
  \qquad
  w_{i+1}\le w_i+\tau
  \quad(0\le i<s-1).
\]
\end{definition}

\begin{definition}[Area, dinv, and defect]
For nonnegative integers \(a,b\), define
\[
  d_\tau(a,b)=
  \begin{cases}
    \max(0,a+\tau-b),&a\le b,\\[2mm]
    \max(0,b+\tau+1-a),&a>b.
  \end{cases}
\]
For a length-\(s\) word \(W\), put
\[
  \area(W)=\sum_{j=0}^{s-1}w_j,
  \qquad
  \dinv_\tau(W)=\sum_{0\le i<j<s}d_\tau(w_i,w_j),
\]
\[
  M_s=\tau\binom{s}{2},
  \qquad
  \defc_\tau(W)=M_s-\area(W)-\dinv_\tau(W).
\]
For a fixed defect value \(d\), the middle area is
\[
  L_d=\left\lfloor\frac{M_s-d}{2}\right\rfloor.
\]
A normalized word \(D\) of defect \(d\) is \emph{strictly below the
middle} when \(\area(D)<L_d\).
\end{definition}

The complementary pair weight is
\[
  K_\tau(a,b)=\tau-d_\tau(a,b)
  =
  \begin{cases}
    \min(\tau,b-a),&a\le b,\\[2mm]
    \min(\tau,a-b-1),&a>b.
  \end{cases}
\]
Thus \(0\le K_\tau(a,b)\le\tau\).  For
\(W=(w_0,\ldots,w_{s-1})\), define the lambda increment
\[
  \lambda_j(W)=\sum_{i<j}K_\tau(w_i,w_j)-w_j.
\]
Since \(K_\tau=\tau-d_\tau\), every word satisfies the exact identity
\begin{equation}
  \defc_\tau(W)=\sum_{j=0}^{s-1}\lambda_j(W).
  \label{eq:defect-lambda}
\end{equation}
It will also be convenient to use the midpoint functional
\[
  J_\tau(W)=2\area(W)+\defc_\tau(W).
\]
Thus \(J_\tau(W)>M_s\) is equivalent to
\(\area(W)>(M_s-\defc_\tau(W))/2\).

\begin{definition}[Records]
An occurrence \(w_j\) is a \emph{record} if \(j=0\), or if
\[
  w_j>\max\{w_0,\ldots,w_{j-1}\}.
\]
Every other occurrence is a \emph{nonrecord}.  The record values of a
normalized word form a chain
\[
  0=r_0<r_1<\cdots<r_{I-1}=m
\]
with \(1\le r_i-r_{i-1}\le\tau\).  The number \(m\) is the
\emph{record peak}.
\end{definition}

Indeed, when a new record is created, the preceding entry is no larger
than the previous record peak, while normalization permits an increase of
at most \(\tau\).

\begin{definition}[The local tests \(\East_3\) and \(\East_5\)]
On a three-window \((x_{-1},x_0,x_1)\), the test \(\East_3\) is defined
exactly when
\[
  x_0\le x_1+\tau.
\]
When defined, it leaves the window unchanged.

Suppose \(\East_3\) has failed on the inner three entries of a five-window
\[
  (a,b,c,p,q),
  \qquad c>p+\tau.
\]
Define
\[
  \bk_\tau(b,c)=
  \begin{cases}
    (c,b),&b>c+\tau,\\
    (b,c),&b\le c+\tau.
  \end{cases}
\]
The map \(\East_5\) is defined as follows.
\begin{enumerate}[label=\textup{(\alph*)}]
  \item If \(b\le p+\tau\), then \(\East_5\) is defined exactly when
  \(b\le q+\tau\), and its output is
  \[
    (a,p,c,b,q).
  \]
  \item If \(b>p+\tau\), write
  \((b',c')=\bk_\tau(b,c)\).  Then \(\East_5\) is defined exactly when
  \(c'\le q+\tau\), and its output is
  \[
    (a,p,b',c',q).
  \]
\end{enumerate}
If the required final inequality fails, \(\East_5\) is said to
\emph{fail}.
\end{definition}

\section{Kernel and record lemmas}

\begin{lemma}[A record chain pays a target]
\label{lem:record-payment}
Let
\[
  0=r_0<r_1<\cdots<r_{I-1}=m,
  \qquad r_i-r_{i-1}\le\tau.
\]
For \(0\le u\le m+\tau\), set
\[
  A(u)=\sum_{r_i\le u}K_\tau(r_i,u),
  \qquad
  G(u)=\sum_{i=0}^{I-1}K_\tau(r_i,u).
\]
Then
\begin{align}
  A(u)&\ge u,
  \label{eq:record-basic}\\
  G(u)&\ge u+(m-1-u)_+.
  \label{eq:record-strong}
\end{align}
\end{lemma}

\begin{proof}
If \(u\ge m\), then for \(i<I-1\),
\[
  K_\tau(r_i,u)=\min(\tau,u-r_i)\ge r_{i+1}-r_i,
\]
and \(K_\tau(m,u)=u-m\).  Summing gives \(A(u)=G(u)\ge u\).

Now suppose \(u<m\), and let \(h\) be the largest index with
\(r_h\le u\).  For \(i<h\),
\[
  K_\tau(r_i,u)\ge r_{i+1}-r_i,
\]
while \(K_\tau(r_h,u)=u-r_h\).  Hence the records at or below \(u\)
pay at least \(u\), proving \eqref{eq:record-basic}.

It remains to prove the extra term in \eqref{eq:record-strong}.  This term
is zero when \(u\ge m-1\).  If \(u\le m-2\), then
\[
  K_\tau(r_{h+1},u)=r_{h+1}-u-1,
\]
because \(r_{h+1}\le r_h+\tau\le u+\tau\).  For every
\(i\ge h+2\),
\[
  K_\tau(r_i,u)\ge r_i-r_{i-1}.
\]
Thus the records above \(u\) contribute at least
\[
  r_{h+1}-u-1+(m-r_{h+1})=m-u-1.
\]
Together with \eqref{eq:record-basic}, this proves
\eqref{eq:record-strong}.
\end{proof}

\begin{corollary}[Lambda nonnegativity]
\label{cor:lambda-nonnegative}
If \(W\) is normalized, then \(\lambda_j(W)\ge0\) at every position.
In particular, \(\defc_\tau(W)\ge0\).
\end{corollary}

\begin{proof}
The assertion is clear at the first position.  At a later position, let
\(m\) be the preceding record peak and put \(u=w_j\).  If \(u\) is a new
record, normalization gives \(u\le m+\tau\); otherwise \(u\le m\).
The preceding records at or below \(u\) contribute at least \(u\) by
\eqref{eq:record-basic}, which cancels the subtraction in
\(\lambda_j(W)\).
\end{proof}

\begin{lemma}[Forward and reverse far-pair bounds]
\label{lem:far-pairs}
If \(H\ge L+\tau+1\), then for every nonnegative integer \(x\),
\begin{align}
  K_\tau(x,H)+K_\tau(x,L)
  &\ge \tau+\min\bigl(\tau,(L-x)_+\bigr),
  \label{eq:forward-far}\\
  K_\tau(H,x)+K_\tau(L,x)&\ge\tau.
  \label{eq:reverse-far}
\end{align}
\end{lemma}

\begin{proof}
If \(x\le L\), then \(K_\tau(x,H)=\tau\), and
\eqref{eq:forward-far} is an equality.  If \(L<x\le H\), put
\(a=x-L-1\) and \(b=H-x\).  Then \(a,b\ge0\) and
\(a+b=H-L-1\ge\tau\), so
\[
  K_\tau(x,L)+K_\tau(x,H)
  =\min(\tau,a)+\min(\tau,b)\ge\tau.
\]
If \(x>H\), then \(K_\tau(x,L)=\tau\).  This proves
\eqref{eq:forward-far}.

For \eqref{eq:reverse-far}, the result is immediate when \(x<L\) or
\(x\ge H\), since one of the two terms is then \(\tau\).  If
\(L\le x<H\), the two uncapped quantities are \(x-L\) and
\(H-x-1\), whose sum is \(H-L-1\ge\tau\).  Capping each at \(\tau\)
therefore leaves a sum of at least \(\tau\).
\end{proof}

\begin{lemma}[Consequences of an \(\East_5\) failure]
\label{lem:east5-consequences}
Suppose
\[
  c>p+\tau,
  \qquad c\le b+\tau,
\]
and \(\East_5\) fails on \((a,b,c,p,q)\).  Then:
\begin{enumerate}[label=\textup{(\roman*)}]
  \item the four entries \(b,c,p,q\) can be paired as
  \[
    (H_1,L_1),\qquad(H_2,L_2),
  \]
  where \(\{H_1,H_2\}=\{b,c\}\),
  \(\{L_1,L_2\}=\{p,q\}\), and
  \[
    H_i\ge L_i+\tau+1\qquad(i=1,2);
  \]
  \item if
  \[
    U=K_\tau(b,c)+K_\tau(b,p)+K_\tau(b,q)
      +K_\tau(c,p)+K_\tau(c,q)+K_\tau(p,q),
  \]
  then
  \begin{equation}
    U\ge4\tau+(c-b-1)_+.
    \label{eq:internal-U}
  \end{equation}
\end{enumerate}
\end{lemma}

\begin{proof}
There are three failure modes.

If \(b\le p+\tau\), failure means \(b>q+\tau\).  The far pairs are
\((c,p)\) and \((b,q)\).  Since \(c>p+\tau\) and
\(c\le b+\tau\), one has
\[
  1\le r:=b-p\le\tau,
  \qquad c>b,
  \qquad q\le b-\tau-1\le p-1.
\]
The three terms \(K_\tau(c,p),K_\tau(b,q),K_\tau(c,q)\) equal
\(\tau\), while
\[
  K_\tau(b,c)=c-b,
  \qquad
  K_\tau(b,p)=r-1,
  \qquad
  K_\tau(p,q)\ge\tau-r.
\]
Thus \(U\ge4\tau+c-b-1\).

If \(b>p+\tau\) and \(b\le c+\tau\), failure means
\(c>q+\tau\).  The far pairs are \((b,p)\) and \((c,q)\).
Moreover,
\[
  K_\tau(b,p)=K_\tau(c,p)=K_\tau(c,q)=\tau.
\]
Applying \eqref{eq:reverse-far} to the far pair \((b,p)\) with target
\(q\) gives
\[
  K_\tau(b,q)+K_\tau(p,q)\ge\tau.
\]
Finally, normalization gives
\[
  K_\tau(b,c)\ge(c-b-1)_+.
\]
Hence \eqref{eq:internal-U} holds.

If \(b>p+\tau\) and \(b>c+\tau\), failure means
\(b>q+\tau\).  The far pairs are \((c,p)\) and \((b,q)\), and the
four terms
\[
  K_\tau(b,c),\quad K_\tau(b,p),\quad
  K_\tau(b,q),\quad K_\tau(c,p)
\]
all equal \(\tau\).  Here \((c-b-1)_+=0\), so
\eqref{eq:internal-U} again follows.
\end{proof}

\section{Suffix and prefix ledgers}

\begin{lemma}[Suffix ledger for a failed window]
\label{lem:suffix-ledger}
Let
\[
  \Sigma=P:b:c:p:q
\]
be a length-\(s\) word, where \(n=|P|=s-4\ge1\), \(P:b:c\) is
normalized, \(c>p+\tau\), and \(\East_5\) fails on the final five
entries.  Let \(m\) be the record peak of \(P\), and let \(N\) be the
number of nonrecords in \(P\).  Define
\[
  S=b+c+p+q,
\]
\[
  E_u=\sum_{x\text{ occurring in }P}
       \min\bigl(\tau,(u-x)_+\bigr),
  \qquad E=E_p+E_q.
\]
\[
  X=(b+c-2m-2\tau-1)_+,
\]
\[
  Y=(m-1-b)_++(m-1-p)_++(m-1-q)_+.
\]
Let \(\Lambda\) be the sum of the final four lambda increments of
\(\Sigma\).  Then
\begin{align}
  E_p&\ge p,
  &E_q&\ge q,
  \label{eq:E-pays}\\
  \Lambda&\ge2\tau n+4\tau+E+X-S,
  \label{eq:suffix-refined}\\
  \Lambda&\ge2\tau N+3\tau+Y.
  \label{eq:suffix-low}
\end{align}
\end{lemma}

\begin{proof}
Let \(0=r_0<\cdots<r_{I-1}=m\) be the record chain of \(P\).  If
\(a\) is the final entry of \(P\), normalization gives
\[
  b\le a+\tau\le m+\tau,
  \qquad c\le b+\tau\le m+2\tau.
\]
Since \(c>p+\tau\),
\[
  p\le m+\tau-1.
\]
By Lemma~\ref{lem:east5-consequences}, \(q\) is the lower member of a
far pair whose upper member is \(b\) or \(c\); hence also
\[
  q\le m+\tau-1.
\]
The records at or below \(p\), respectively \(q\), contribute to
\(E_p\), respectively \(E_q\).  Therefore
\eqref{eq:record-basic} proves \eqref{eq:E-pays}.

For each occurrence \(x\) in \(P\), pair the four suffix entries as in
Lemma~\ref{lem:east5-consequences}.  Applying
\eqref{eq:forward-far} to the two pairs gives
\[
  K_\tau(x,b)+K_\tau(x,c)+K_\tau(x,p)+K_\tau(x,q)
  \ge2\tau+
  \min\bigl(\tau,(p-x)_+\bigr)+
  \min\bigl(\tau,(q-x)_+\bigr).
\]
Summing over \(P\) gives \(2\tau n+E\).  The six internal terms among
\(b,c,p,q\) contribute \(U\).  Moreover, since \(b\le m+\tau\),
\[
  b+c-2m-2\tau-1\le c-b-1,
\]
so \(X\le(c-b-1)_+\).  Equation \eqref{eq:internal-U} now gives
\[
  \Lambda\ge2\tau n+E+4\tau+X-S,
\]
which is \eqref{eq:suffix-refined}.

For \eqref{eq:suffix-low}, write
\[
  G(u)=\sum_{i=0}^{I-1}K_\tau(r_i,u).
\]
By \eqref{eq:record-strong},
\begin{equation}
  G(b)+G(p)+G(q)\ge b+p+q+Y.
  \label{eq:G-bpq}
\end{equation}
Every nonrecord occurrence of \(P\) contributes at least \(2\tau\) to
the four suffix lambda sums, by the two far-pair bounds.

It remains to account for the subtraction of \(c\).  We claim that
\begin{equation}
  G(c)+U\ge c+3\tau.
  \label{eq:c-payment}
\end{equation}
First,
\[
  G(c)+K_\tau(b,c)\ge c.
\]
Indeed, if \(b\le m\), then \(c\le m+\tau\), so \(G(c)\ge c\) by
\eqref{eq:record-basic}.  If \(b>m\), append \(b\) to the record chain;
since \(b\le m+\tau\) and \(c\le b+\tau\), the extended chain pays
\(c\).  Therefore, using \(U\ge4\tau\) and
\(K_\tau(b,c)\le\tau\),
\[
  G(c)+U
  =\bigl(G(c)+K_\tau(b,c)\bigr)
   +\bigl(U-K_\tau(b,c)\bigr)
  \ge c+3\tau.
\]

Combining \eqref{eq:G-bpq}, the \(2\tau N\) contribution of the
nonrecords, and \eqref{eq:c-payment}, and then subtracting \(S\), yields
\eqref{eq:suffix-low}.
\end{proof}

\begin{lemma}[Peak and lower-target mass]
\label{lem:peak-and-E}
In the setting of Lemma~\ref{lem:suffix-ledger}, let
\[
  \Delta=\sum_{j\text{ in }P}\lambda_j(P),
  \qquad d=\defc_\tau(\Sigma)=\Delta+\Lambda.
\]
If
\begin{equation}
  d\le\tau(n+2)-1,
  \label{eq:defect-range-n}
\end{equation}
then
\begin{align}
  m&\ge\frac{\tau n}{2},
  \label{eq:peak-bound}\\
  E&\ge\tau(1+2N)-1+\Delta.
  \label{eq:E-bound}
\end{align}
\end{lemma}

\begin{proof}
The prefix \(P\) is normalized, so \(\Delta\ge0\) by
Corollary~\ref{cor:lambda-nonnegative}.  From
\eqref{eq:suffix-refined} and \eqref{eq:E-pays},
\[
  d=\Delta+\Lambda
  \ge2\tau n+4\tau+X-(b+c).
\]
Using \eqref{eq:defect-range-n},
\begin{equation}
  b+c-X\ge\tau(n+2)+1.
  \label{eq:bc-lower}
\end{equation}
By the definition of \(X\),
\[
  b+c-X\le2m+2\tau+1.
\]
Combining this with \eqref{eq:bc-lower} proves
\eqref{eq:peak-bound}.

Next, \eqref{eq:suffix-low} and \eqref{eq:defect-range-n} give
\begin{equation}
  Y\le\tau(n-1-2N)-1-\Delta.
  \label{eq:Y-upper}
\end{equation}
For \(u=p,q\), equation \eqref{eq:E-pays} implies
\[
  E_u+(m-1-u)_+\ge m-1.
\]
Since the two displayed positive parts are contained in \(Y\),
\[
  E\ge2m-2-Y.
\]
Applying \eqref{eq:peak-bound} and \eqref{eq:Y-upper} gives
\[
  E\ge\tau n-2-\bigl(\tau(n-1-2N)-1-\Delta\bigr)
   =\tau(1+2N)-1+\Delta.
\]
\end{proof}

\begin{lemma}[Prefix cancellation]
\label{lem:prefix-cancellation}
Let \(P\) be a normalized word with record values
\[
  0=r_0<r_1<\cdots<r_{I-1}=m
\]
and \(N\) nonrecord occurrences.  Put
\[
  T=\sum_{0\le i<j<I}K_\tau(r_i,r_j),
  \qquad
  \Delta=\sum_{j\text{ in }P}\lambda_j(P).
\]
Then
\begin{equation}
  \area(P)+\Delta\ge T+N(m-1).
  \label{eq:prefix-cancellation}
\end{equation}
\end{lemma}

\begin{proof}
Write \(P=(w_0,\ldots,w_{n-1})\).  By the definition of the lambda
increments,
\[
  \area(P)+\Delta
  =\sum_{i<j}K_\tau(w_i,w_j),
\]
the total \(K_\tau\)-weight of all chronologically ordered pairs in
\(P\).

Fix a nonrecord occurrence of value \(z\), and let \(r_t\) be the latest
record preceding it.  Its total \(K_\tau\)-weight with the records, in
the correct chronological orientation, is
\[
  C(z)=
  \sum_{i=0}^{t}K_\tau(r_i,z)
  +\sum_{i=t+1}^{I-1}K_\tau(z,r_i).
\]
We claim
\begin{equation}
  C(z)\ge m-1.
  \label{eq:cross-payment}
\end{equation}
If \(z=r_t\), the preceding records pay \(z\) by
\eqref{eq:record-basic}, and the later records contribute at least
\(m-r_t\) by telescoping the record gaps.  Hence \(C(z)\ge m\).

If \(z<r_t\), let \(h\) be the largest index with \(r_h\le z\).  The
records through \(r_h\) pay \(z\).  The records from \(r_{h+1}\) through
\(r_t\) contribute at least \(r_t-z-1\): the first term is
\(r_{h+1}-z-1\), and each later term dominates the next record gap.
The records after the occurrence of \(z\) contribute at least
\(m-r_t\).  Thus \(C(z)\ge m-1\), proving
\eqref{eq:cross-payment}.

Summing \(C(z)\) over the \(N\) nonrecords gives exactly the
record-to-nonrecord plus nonrecord-to-record \(K_\tau\)-weight.  If
\(K_{N\to N}\) denotes the nonrecord-to-nonrecord weight, then
\[
  \area(P)+\Delta
  =T+\sum_{z\text{ nonrecord}}C(z)+K_{N\to N}
  \ge T+N(m-1),
\]
since \(K_{N\to N}\ge0\).
\end{proof}

\begin{lemma}[Record-energy bound]
\label{lem:record-energy}
Let
\[
  0=r_0<r_1<\cdots<r_{I-1}=m,
  \qquad r_i-r_{i-1}\le\tau,
\]
and set \(x=m/\tau\).  Then
\begin{equation}
  \sum_{0\le i<j<I}K_\tau(r_i,r_j)
  \ge
  \tau\left(Ix-\frac{x(x+1)}2-\frac18\right).
  \label{eq:record-energy}
\end{equation}
\end{lemma}

\begin{proof}
Put
\[
  u_j=\frac{r_j-r_{j-1}}{\tau}
  \qquad(1\le j<I).
\]
Then \(0<u_j\le1\), \(\sum_j u_j=x\), and the left side of
\eqref{eq:record-energy}, divided by \(\tau\), is
\[
  \Phi(u_1,\ldots,u_{I-1})
  =\sum_{0\le i<j<I}
   \min\bigl(1,u_{i+1}+\cdots+u_j\bigr).
\]
Each summand is \(t\mapsto\min(1,t)\) applied to a linear form, so
\(\Phi\) is concave.  Relax the domain to
\[
  0\le u_j\le1,
  \qquad \sum_j u_j=x.
\]
A minimum can be chosen at an extreme point: if a minimizer is a
nontrivial convex combination, concavity makes one endpoint another
minimizer.  Such a point has at most one fractional coordinate, since
otherwise a small amount can be transferred between two fractional
coordinates in both directions.

First suppose \(x=k\) is an integer.  Place \(I\) abstract slots
\(0,1,\ldots,I-1\), with \(u_j\) between slots \(j-1\) and \(j\).  The
\(k\) ones split the slots into \(k+1\) nonempty consecutive blocks.  A
summand is one exactly when its endpoint slots lie in different blocks,
so \(\Phi\) counts such pairs.  The number of same-block pairs is at most
\(\binom{I-k}{2}\), attained by block sizes \(I-k,1,\ldots,1\).  Hence
\[
  \Phi\ge q_I(k):=\binom I2-\binom{I-k}{2}
  =Ik-\frac{k(k+1)}2.
\]

Now let \(x=k+\theta\), where \(k=\lfloor x\rfloor\) and
\(0<\theta<1\).  An extreme point has \(k\) ones, one coordinate
\(\theta\), and the rest zero.  Replacing \(\theta\) by zero and one
gives \(v,w\) with \(u=(1-\theta)v+\theta w\).  Therefore
\[
\begin{aligned}
  \Phi(u)
  &\ge(1-\theta)q_I(k)+\theta q_I(k+1)\\
  &=q_I(k)+\theta(I-k-1)\\
  &=q_I(x)-\frac{\theta(1-\theta)}2
  \ge q_I(x)-\frac18,
\end{aligned}
\]
where \(q_I(t)=It-t(t+1)/2\).  This is \eqref{eq:record-energy}.
\end{proof}

\section{The uniform local obstruction}

\begin{theorem}[A failed \(\East_5\) window lies above the midpoint]
\label{thm:local-obstruction}
Let \(\tau\ge2\) and \(s\ge5\).  Suppose
\[
  \Sigma=P:b:c:p:q
\]
is a length-\(s\) word such that
\begin{enumerate}[label=\textup{(\roman*)}]
  \item \(P:b:c\) is a normalized step-\(\tau\) Dyck word;
  \item \(p,q\) are nonnegative integers;
  \item \(c>p+\tau\);
  \item \(\East_5\) fails on the final five entries;
  \item \(d=\defc_\tau(\Sigma)\le\tau(s-2)-1\).
\end{enumerate}
Then
\begin{equation}
  J_\tau(\Sigma)=2\area(\Sigma)+d>M_s=\tau\binom{s}{2}.
  \label{eq:local-conclusion}
\end{equation}
Equivalently,
\[
  \area(\Sigma)>\frac{M_s-d}{2}.
\]
\end{theorem}

\begin{proof}
Put \(n=s-4\), so
\[
  d\le\tau(n+2)-1.
\]
Let \(m\) be the record peak of \(P\), let \(N\) be its number of
nonrecords, and put \(I=n-N\) and \(x=m/\tau\).  Let \(\Delta\) be the
prefix lambda sum, \(\Lambda=d-\Delta\) the final-four lambda sum, \(E\)
the upward prefix-to-\(p,q\) sum, \(X=(b+c-2m-2\tau-1)_+\), and \(T\)
the record--record \(K_\tau\)-sum.

If \(n=1\), then \(P=(0)\), so \(N=0\) and \(\Delta=0\).  Equation
\eqref{eq:suffix-low} gives \(d=\Lambda\ge3\tau\), contradicting
\(d\le3\tau-1\).  Hence no failed window exists in this case.

Assume \(n\ge2\).  By Lemma~\ref{lem:peak-and-E},
\begin{equation}
  x\ge\frac n2,
  \qquad
  E\ge\tau(1+2N)-1+\Delta.
  \label{eq:x-E-summary}
\end{equation}
Since the \(I-1\) record gaps are at most \(\tau\),
\begin{equation}
  x\le I-1\le n-1.
  \label{eq:x-upper}
\end{equation}

Rearranging \eqref{eq:suffix-refined} and using
\(d=\Delta+\Lambda\) gives
\[
  b+c+p+q\ge2\tau(n+2)+E+X+\Delta-d.
\]
Adding this to \eqref{eq:prefix-cancellation} yields
\begin{equation}
  \area(\Sigma)
  \ge T+N(m-1)+2\tau(n+2)+E+X-d.
  \label{eq:area-master}
\end{equation}
Using \(X\ge0\) and \(-d\ge-\tau(n+2)+1\),
\begin{align}
  J_\tau(\Sigma)-M_s
  &\ge2T+2N(m-1)+4\tau(n+2)+2E+2X-d-M_s\notag\\
  &\ge2T+2N(m-1)+3\tau(n+2)+1+2E-M_s.
  \label{eq:J-master}
\end{align}

By Lemma~\ref{lem:record-energy},
\(2T\ge\tau(2Ix-x(x+1)-1/4)\).  Also
\(2N(m-1)=2\tau Nx-2N\), \(I+N=n\), and
\(M_s=\tau(n+4)(n+3)/2\).  Thus \eqref{eq:J-master} becomes
\begin{align}
  J_\tau(\Sigma)-M_s
  \ge\tau\left(
    2nx-x(x+1)-\frac14+3(n+2)
    -\frac{(n+4)(n+3)}2
  \right)
  +1-2N+2E.
  \label{eq:J-x}
\end{align}
For \(f(x)=2nx-x(x+1)\), \eqref{eq:x-upper} gives
\(f'(x)=2n-2x-1\ge1\).  Hence \eqref{eq:x-E-summary} gives
\(f(x)\ge f(n/2)=3n^2/4-n/2\).  Substitution into \eqref{eq:J-x} gives
\begin{equation}
  J_\tau(\Sigma)-M_s
  \ge
  \frac{\tau(n^2-4n-1)}4+1-2N+2E.
  \label{eq:J-before-E}
\end{equation}
Finally, \eqref{eq:x-E-summary} and \(\Delta\ge0\) give
\(E\ge\tau(1+2N)-1\).  Substitution into
\eqref{eq:J-before-E} yields
\begin{align*}
  J_\tau(\Sigma)-M_s
  &\ge
  \frac{\tau\bigl((n-2)^2+3\bigr)}4
  +(4\tau-2)N-1\\
  &\ge \frac{3\tau}{4}-1
  \ge\frac12>0.
\end{align*}
This proves \eqref{eq:local-conclusion}.
\end{proof}

\end{document}
